\input{../tex-inputs/latex-leseansicht-vorspann}

\section[Arthur Schnitzler an Gustav Schwarzkopf, 12. 10. 1905]{L04020 Arthur Schnitzler an Gustav Schwarzkopf, 12. 10. 1905}
\nopagebreak\mylabel{L04020v}
\rehead{ }\normalsize\beginnumbering\briefempfaengerindex{Schwarzkopf, Gustav@\textsc{Schwarzkopf, Gustav}!zzzSchnitzler, Arthur@\emph{von Arthur Schnitzler}!1905-10-121@{12. 10. 1905}|(be}
\toendnotes[C]{\smallbreak\pagebreak[2]}
\correspDesc{Versand  durch Arthur Schnitzler am 12. 10. 1905 in Wien
\newline{}Erhalt  durch Gustav Schwarzkopf im Zeitraum [12. 10. 1905 – 15. 10. 1905?] in Wien}\toendnotes[C]{\smallbreak}
\Standort{CUL, Schnitzler, B 96.}
\physDesc{Briefkarte, 233 Zeichen
\newline{}Handschrift: schwarze Tinte, deutsche Kurrent}\toendnotes[C]{\smallbreak}
\pstart
           {\pb}\textcolor{gray}{\textbf{Dr. Arthur Schnitzler}}\hfill 12. X. 1905\pend
           
\pstart
           \textcolor{gray}{\textbf{Wien,
                        XVIII. Spoettelgasse 7\oindex{Wien@\textbf{Wien}!XVIII., Währing@\textbf{XVIII., Währing}!Edmund-Weiß-Gasse@\textbf{Edmund-Weiß-Gasse}, \emph{Straße}|pw}.}}\pend
           \vspace{0.5em}
\pstart
           lieber Guſtav, beifolgd 2 Sitze zur Kakadupremière\pwindex{Schnitzler, Arthur 15.\,5.\,1862 Wien – 21.\,10.\,1931 ebd.@\textsc{Schnitzler, Arthur} (15.\,5.\,1862 Wien – 21.\,10.\,1931 ebd.), \emph{Schriftsteller, Mediziner}!grüne Kakadu. Groteske in einem Akt@\strich\emph{Der grüne Kakadu. Groteske in einem Akt}|pw}\eventindex{Volkstheater@\textbf{Volkstheater}!Premiere von Der zerbrochene Krug, Der grüne Kakadu, 14.10.1905@Premiere von Der zerbrochene Krug, Der grüne Kakadu, 14.10.1905|pwv}. – Ob ich morgen den Sitz
               zum Zwiſchenſpiel\pwindex{Schnitzler, Arthur 15.\,5.\,1862 Wien – 21.\,10.\,1931 ebd.@\textsc{Schnitzler, Arthur} (15.\,5.\,1862 Wien – 21.\,10.\,1931 ebd.), \emph{Schriftsteller, Mediziner}!Zwischenspiel. Komödie in drei Akten@\strich\emph{Zwischenspiel. Komödie in drei Akten}|pw}\eventindex{Burgtheater@\textbf{Burgtheater}!2. Aufführung von Zwischenspiel, 13.10.1905@2. Aufführung von Zwischenspiel, 13.10.1905|pw} an Ihren \label{K_L04020-1v}\edtext{Bruder\pwindex{Schwarzkopf, Max 12.\,6.\,1857 Wien – 14.\,4.\,1928 ebd.@\textsc{Schwarzkopf, Max} (12.\,6.\,1857 Wien – 14.\,4.\,1928 ebd.), \emph{Rechtsanwalt}|pwv}{ }ſenden ka{\geminationn}}{\lemma{\textnormal{\emph{Bruder senden kann}}}\Cendnote{\textnormal{Siehe XXXX Auszeichnungsfehler: Dokument L04017 nicht gefunden.}}}\label{K_L04020-1}, weiſs ich in
               diesem Moment noch nicht ganz{ }ſicher; we{\geminationn} nicht diesmal,
               nächſte Woche beſti{\geminationm}t.\pend
           \pstart Herzlich Ihr \spacefill\mbox{A.}\pend{}\selectlanguage{ngerman}\endnumbering\briefempfaengerindex{Schwarzkopf, Gustav@\textsc{Schwarzkopf, Gustav}!zzzSchnitzler, Arthur@\emph{von Arthur Schnitzler}!1905-10-121@{12. 10. 1905}|)be}\mylabel{L04020h}
\begin{anhang}
\end{anhang}\newcommand{\dateiname}{L04020}\newcommand{\titel}{Arthur Schnitzler an Gustav Schwarzkopf, 12. 10. 1905}\newcommand{\editorInnen}{Herausgegeben von Jahnke, SelmaMüller, Martin Anton}\input{../tex-inputs/latex-leseansicht-abspann}
